\section*{الفصل \ref{ch:b3:genetic-engineering} --- الهندسة الوراثية والتقانة الحيوية}

\begin{solution}{exo:b3:genetic-engineering:1}
منشأ تضاعف، حتى ينسخ المضيف \omterm{def:b3:genetic-engineering:tools}{البلازميد} ويمرره
إلى الخلايا البنات؛ \omterm{def:b3:genetic-engineering:tools}{وواسم انتقائي}، وهو عادةً مقاومة لمضاد
حيوي، حتى لا ينمو إلا ما يحمل \omterm{def:b3:genetic-engineering:tools}{البلازميد}؛ وموقع قطع
فريد (أو عنقود مواقع) تُوصل عنده المقحَمة
بلا قطع \omterm{def:b3:genetic-engineering:tools}{للبلازميد} في مكان آخر.
\end{solution}

\begin{solution}{exo:b3:genetic-engineering:2}
المتتالية المعينة من ست قواعد ترد مرة كل $4^{6} = 4096$ زوج قواعد في
DNA عشوائي: أي نحو $4.6\times 10^{6}/4096 \approx 1100$ قطعة متوسطها
\qty{4.1}{kb}. والجينومات الحقيقية ليست عشوائية --- فتركيب القواعد
وانحياز الرامزات وتجنب بعض المتتاليات تزيح العدد (وعدد
مواقع EcoRI الفعلي في \emph{E.~coli} بضع مئات).
\end{solution}

\begin{solution}{exo:b3:genetic-engineering:3}
التمسيخ عند \qty{95}{\celsius} يفصل الطاقين؛ والتلاحم عند
\qtyrange{50}{65}{\celsius} يتيح للبادئتين أن تزدوجا بموقعيهما؛
والاستطالة عند \qty{72}{\celsius} تتيح للبوليميراز أن ينسخ من كل
\omterm{def:b3:genetic-engineering:pcr}{بادئة}. والبوليميراز العادي كان ليُهدم عند \qty{95}{\celsius}
في الدورة الأولى ولوجب إضافته من جديد ثلاثين مرة؛ أما Taq
فينجو من الدورات كلها.
\end{solution}

\begin{solution}{exo:b3:genetic-engineering:4}
\omterm{def:b3:genetic-engineering:crispr}{RNA دليل} قواعده العشرون تطابق الهدف، وموقع \omterm{def:b3:genetic-engineering:crispr}{PAM} (من نمط NGG)
يلي الهدف مباشرةً عند 3$'$ على الطاق غير الهدف. وبعد
\omterm{def:b3:dna-repair:lesions}{الكسر المزدوج الطاق}: إما وصل النهايات، فيترك إقحامًا أو
حذفًا صغيرًا ويهدم عادةً إطار قراءة المورّثة؛ وإما \omterm{def:b3:dna-repair:dsb}{إعادة تركيب
متماثلة} مع قالب يحمل التسلسل المرغوب، فتكتبه
فيه.
\end{solution}

\begin{solution}{exo:b3:genetic-engineering:5}
$50\times 1.9^{35} = 50\times 5.7\times 10^{9} \approx 3\times 10^{11}$
نسخة. و$\Delta C_{T} = 7.3$؛ والنسبة $1.95^{7.3} = e^{7.3\ln 1.95}
\approx 130$: أي كان في العينة الأولى نحو $130$ ضعف القالب.
\end{solution}

\begin{solution}{exo:b3:genetic-engineering:6}
البكتيريا لا تستطيع التشذيب: فالنسخة الجينومية بإنتروناتها
تُستنسخ رسالةً غير صالحة، فيُستعمل cDNA الخالي من الإنترونات، المنسوخ
عن mRNA الناضج. ومن العقبات الأخرى: أن بوليميراز البكتيريا لا
يتعرف المروّج البشري (فزوّده بمروّج بكتيري وموقع
ربط ريبوزومي)؛ وأن انحياز الرامزات البشري مختلف (فركّب المورّثة
برامزات المضيف المفضلة)؛ وأن البروتين قد لا ينطوي أو قد
يتكتل (فاخفض الحرارة، أو الحمه بشريك ذواب، أو أعد طيّه
من أجسام الاحتباس)؛ وأن الغلكزة غائبة (فاستعمل الخميرة أو خلايا
الثدييات إن كانت مهمة)؛ والروابط ثنائية الكبريت تحتاج إلى المحيط المؤكسد حول البلازما.
\end{solution}

\begin{solution}{exo:b3:genetic-engineering:7}
تُحقن الخلايا الجذعية المستهدَفة في كيسة أريمية مضيفة فتختلط
بخلاياها، فيُبنى الفأر من نمطين وراثيين --- أي
خيمري --- ولا ينقل الأليل إلا إذا انحدرت بعض خلاياه الجنسية من
الخلايا المستهدَفة. وتزويج الخيمري بفأر
بري يعطي متغايري لواقح (من نصف أعراسه إذا كان نصف
سلالته الجنسية مستهدَفًا)؛ وتزويج متغايري اللواقح فيما بينهم يعطي ربعًا
متماثل اللواقح محذوف المورّثة.
\end{solution}

\begin{solution}{exo:b3:genetic-engineering:8}
المطابقة التامة: $6.4\times 10^{9}\times 4^{-20}\times (1/16) \approx 4\times
10^{-4}$ موقع (فقاعدتا GG في \omterm{def:b3:genetic-engineering:crispr}{PAM} تكلّفان عامل $16$) --- أي لا شيء
عمليًا. وضمن عدم توافق في موضعين توجد $1 + 60 + 1710 = 1771$
متتالية: أي $1771\times 4\times 10^{-4} \approx 0.7$ موقع مصادف. و
الجينوم الحقيقي ليس عشوائيًا، والدليل يُفحص في مواجهته مباشرةً.
\end{solution}

\begin{solution}{exo:b3:genetic-engineering:9}
$p_{1} = 0.05 + 0.8\times 0.05\times 0.95 = 0.088$؛ و$p_{2} = 0.152$؛
و$p_{3} = 0.255$؛ ثم $0.41$ و$0.60$: أي نحو خمسة أجيال لتجاوز
$0.5$، مقابل التقدير $\ln(1/(2\times 0.05))/\ln 1.8 \approx 4$.
\end{solution}

\begin{solution}{exo:b3:genetic-engineering:10}
وصل النهايات يعمل في كل خلية، ومنها الخلايا الجذعية الساكنة التي
تكون فيها \omterm{def:b3:dna-repair:dsb}{إعادة التركيب المتماثلة}، وهي عملية في S/G2، قليلة الكفاءة؛ وهو
لا يحتاج إلى إيصال قالب وناتجه --- أي تعطيل
المعزِّز أيًا كان --- يؤدي الغرض، في حين يقتضي التصحيح التسلسل
الدقيق ولا يعطي إلا أقلية من الخلايا. كما أن التحرير نفسه ينفع
كل طفرة في الغلوبين $\beta$، منجلية كانت أو ثلاسيمية. والعيب: أنه
يزيل عنصرًا تنظيميًا (بما له من أدوار أخرى)، ويترك
الأليل المنجلي في مكانه، وأن خليط الإقحامات والحذوف غير مضبوط.
\end{solution}

\begin{solution}{exo:b3:genetic-engineering:11}
عند لياقة متماثل اللواقح $1 - s$: $p' = \bigl[p^{2}(1-s) + p(1-p)(1+c)
\bigr]/(1 - sp^{2})$. ومن الندرة تكون متماثلات اللواقح مهملة وينمو
الدافع مثل $(1+c)p$ أيًا كانت $s$؛ أما هل يبلغ التثبيت فيُحسم
قرب $p = 1$، حيث يعود أليل بري نادر تواتره $q$
بمقدار $q' \approx q(1-c)/(1-s)$: فيثبت إذا كان $c > s$، وإلا كان توازنًا
وسطيًا. والمقاومة: كل تحويل فاشل
(وصل نهايات بدل إعادة تركيب) يترك أليلًا لم يعد الدليل
يتعرفه؛ فإذا كانت تلك الأليلات لائقة --- كإقحام أو حذف في الإطار في
منطقة غير جوهرية --- لم تحمل كلفة بينما يحملها الدافع، فيحابيها
الانتقاء وتحل محله. والعلاج أن تُستهدف متتاليات جوهرية تكون
الإقحامات والحذوف فيها مميتة، وأن تُستعمل عدة أدلة دفعةً واحدة.
\end{solution}

\begin{solution}{exo:b3:genetic-engineering:12}
\omterm{def:b3:genetic-engineering:crispr}{Cas9} المحقون في لاقحة قد يعمل بعد الانقسامات الأولى،
فتتلقى خلايا مختلفة إصلاحات مختلفة --- أي فسيفسائية؛ وإصلاح
كل كسر تختاره الخلية، فيعطي وصل النهايات إقحامات وحذوفًا
غير متوقعة وحذوفًا كبيرة أو فقدَ تغاير اللواقح عند القطع؛ وناتج
\omterm{def:b3:dna-repair:dsb}{إعادة التركيب المتماثلة} أقلية؛ وتقع كسور \omterm{prop:b3:genetic-engineering:editors}{خارج الهدف}
في مواضع لا يمكن التنبؤ بها كلها؛ ولا يمكن فحص الجنين
إلا بتسلسل بضع خلايا مأخوذة بالخزعة، وهي لا تنطق عن
البقية. والدليل المطلوب: طرائق تحرّر كل خلية تحريرًا واحدًا
(قبل الطور S الأول) بكفاءات تقارب \qty{100}{\%}،
وتسلسل جينوم كامل لكل خلية من أجنّة حيوانية محرَّرة يبيّن
انعدام أي تغير غير مقصود، ومتابعة طويلة الأمد لحيوانات محرَّرة عبر
الأجيال --- ولا شيء من ذلك موجود.
\end{solution}

\begin{solution}{pb:b3:genetic-engineering:1}
\textbf{1.} مرة كل $4096$ زوج قواعد. و$P(\text{لا موقع}) = (1 - 1/4096)^{1500}
= e^{-0.37} = 0.69$.
\textbf{2.} اختر إنزيمًا لا موقع له في المورّثة؛ أو ضخّم
المورّثة بتفاعل PCR ببادئتين تحملان مواقع قطع جديدة عند نهايتيهما 5$'$
(أو استعمل الاستنساخ الكليل أو المستقل عن الوصل).
\textbf{3.} \qty{0.1}{\micro g}\,$\times 2\times 10^{7} = 2\times
10^{6}$ مستعمرة؛ و\qty{10}{\%} منها، أي $2\times 10^{5}$، تحمل المقحَمة.
\textbf{4.} البيضاء: \qty{10}{\%}. ونصف المقحمات بالاتجاه
الصحيح، ومنه $P(\text{لا واحدة في } n \text{ بيضاء}) = 0.5^{n} \le 0.01$
تعطي $n \ge 6.6$: أي التقط سبعًا.
\textbf{5.} $32$ مضاعفة، أي $2^{32} = 4.3\times 10^{9}$ خلية، و$2.1\times
10^{11}$ \omterm{def:b3:genetic-engineering:tools}{بلازميد}؛ وكل منها $5500\times 650 = 3.6\times 10^{6}$ دالتون $=
5.9\times 10^{-18}$ غرام: أي \qty{1.3}{\micro g} من \omterm{def:b3:genetic-engineering:tools}{البلازميد}.
\textbf{6.} الاستنساخ لا يحتاج إلا إلى التضاعف، ويوفره منشأ
\omterm{def:b3:genetic-engineering:tools}{البلازميد}؛ أما المروّج البشري فلا يقرؤه بوليميراز RNA البكتيري، فلا بد
للتعبير من وضع مروّج بكتيري وموقع ربط ريبوزومي
أمام المتتالية المشفِّرة (الخالية من الإنترونات).
\textbf{7.} $N_{0} = N_{T}/2^{C_{T}}$: أي $10^{12}/2^{28} \approx 3700$
نسخة؛ وعند $C_{T} = 35$، $10^{12}/2^{35} \approx 29$ نسخة.
\textbf{8.} $2^{n} = 10^{12}$: أي $n = 39.9$، فأربعون دورة.
\textbf{9.} $10^{12}/2^{38} \approx 3.6$ نسخة. وبعد نحو $40$
دورة يبلغ العتبةَ جزيءٌ ملوث واحد، أو مصنوعات البادئات،
وتكون العينة التي فيها أقل من نسخة واحدة بلا معنى.
\textbf{10.} الجزيء الواحد يعبر عند الدورة $40$: أي إيجابية. وتُمنع
بغرف ومعدات منفصلة للتهيئة ولمعالجة
النواتج، وبمسار عمل باتجاه واحد، وشواهد سلبية في كل شوط، و
هدم إنزيمي لنواتج التضخيم المنقولة.
\textbf{11.} $3700/0.4 \approx 9000$ نسخة RNA.
\textbf{12.} المبلَّغ عنه $2^{3.3} \approx 10$ أضعاف؛ والحقيقي $1.9^{3.3} =
e^{3.3\ln 1.9} \approx 8.3$ أضعاف.
\textbf{13.} $40\times \qty{35}{\micro g} = \qty{1.4}{mg}$ في اليوم،
و\qty{0.51}{g} في السنة؛ ومع $\times 10^{7}$ مريض: \qty{5100}{kg} في السنة.
\textbf{14.} $5.1\times 10^{6}/5800 \approx 880$ مول؛ و$5.3\times 10^{26}$
جزيء.
\textbf{15.} $5.1\times 10^{6}$ غرام $/ (20$ غرام/لتر$) = 2.6\times 10^{5}$ لتر
$= \qty{260}{m^3}$ --- أي عُشر المسبح.
\textbf{16.} $5100\times 8000 = 4.1\times 10^{7}$ كيلوغرام من البنكرياس، وكل واحد
\qty{0.1}{kg}: أي $4\times 10^{8}$ خنزير في السنة.
\textbf{17.} الهرمون المؤتلف مطابق للبشري، فلا
يثير أجسامًا مضادة ولا يسبب تفاعلات تحسسية؛ ولا يحمل
ممرضات حيوانية (فيروسات أو بريونات)؛ ولا يتوقف توفره على الذبح؛
ويمكن تحسين تسلسله.
\textbf{18.} البقايا التي تلامس المستقبل لم تتغير، فيكون
الارتباط والإشارة كما في الأنسولين؛ أما البقايا المبدَّلة فتقع
على السطح حيث تتجمع جزيئات الأنسولين ثنائيات
وسداسيات، وتبديلها يغيّر سرعة تفكك المستودع المحقون
إلى وحيدات قابلة للامتصاص.
\textbf{19.} $6.4\times 10^{9}\times 4^{-20}/16 \approx 3.6\times
10^{-4}$: أي لا مطابقة تامة مصادفة.
\textbf{20.} $1 + 60 + 1710 + 27\times 1140 = \num{32551}$ متتالية؛
و$\num{32551}\times 3.6\times 10^{-4} \approx 12$ موقعًا مصادفًا \omterm{prop:b3:genetic-engineering:editors}{خارج
الهدف}.
\textbf{21.} $10^{7}\times 10^{-4} = 1000$ خلية. والخلية الجذعية تعيش
وتنقسم طوال عمر المريض؛ والتي فقدت كابحًا
للورم، أو اكتسبت انتقالًا صبغيًا، قد تؤسس نسيلة تصير
ابيضاضًا.
\textbf{22.} $p_{1} = 0.0019$، و$p_{2} = 0.0036$، و$p_{3} = 0.0069$؛
والنمو الهندسي بمقدار $1.9$ في الجيل يعطي $\ln 500/\ln 1.9
\approx 10$؛ وبتكرار العلاقة يتجاوز التواتر $0.5$ عند
الجيل $11$ ($0.45 \to 0.68$).
\textbf{23.} نحو سنة. وكلفة \qty{20}{\%} في متماثلات اللواقح أقل
من $c = 0.9$، فيبلغ الدافع التثبيت مع ذلك، لكن أبطأ، بينما
تهبط لياقة الجمهرة الوسطى --- وهو غرض دافع الكبح؛ أما
كلفة تفوق $c$ فكانت لتوقفه عند تواتر وسطي.
\textbf{24.} $10^{6}\times 0.1 = 10^{5}$ أليل مقاومة في جيل
واحد. وهي غير قابلة للقطع، وإن كانت لائقة فبلا كلفة بينما الدافع
مكلف، فتُحابى وتنتشر ويُقصى الدافع --- ولهذا
تستهدف الدوافع مواقع تكون نواتج وصل النهايات فيها مميتة، وبعدة
أدلة دفعةً واحدة.
\textbf{25.} نحو $3700$ نسخة في العينة الأولى؛ ونحو \qty{260}{m^3}
من المخمّرة في السنة لأنسولين العالم؛ ونحو $11$ جيلًا، أي سنة
بعوض، للدافع.
\end{solution}
